\documentclass[a4paper,11pt]{ltjsarticle}
\usepackage{luatexja}
\usepackage{luatexja-fontspec}
\usepackage{amsmath,amssymb,amsthm}
\usepackage{mathtools}
\usepackage{booktabs}
\usepackage{longtable}
\usepackage{array}
\usepackage{hyperref}
\usepackage{graphicx}
\usepackage{float}
\usepackage{geometry}
\geometry{margin=25mm}

\hypersetup{
  colorlinks=true,
  linkcolor=blue,
  urlcolor=blue,
  citecolor=blue
}

\theoremstyle{definition}
\newtheorem{theorem}{定理}[section]
\newtheorem{proposition}[theorem]{命題}
\newtheorem{lemma}[theorem]{補題}
\newtheorem{corollary}[theorem]{系}
\newtheorem{definition}[theorem]{定義}
\newtheorem{remark}[theorem]{注意}
\newtheorem{observation}[theorem]{観察}

\setlength{\parindent}{1em}
\setlength{\parskip}{0.5em}

\begin{document}

\section{\texorpdfstring{離散中心投影における第5次元の幾何学的役割：電荷の量子化と \(1/r^2\) 法則}{離散中心投影における第5次元の幾何学的役割：電荷の量子化と 1/r\^{}2 法則}}\label{ux96e2ux6563ux4e2dux5fc3ux6295ux5f71ux306bux304aux3051ux308bux7b2c5ux6b21ux5143ux306eux5e7eux4f55ux5b66ux7684ux5f79ux5272ux96fbux8377ux306eux91cfux5b50ux5316ux3068-1r2-ux6cd5ux5247}

\textbf{著者：} Noriaki Kihara（木原 範昭）\\
\textbf{所属：} WF System Co., Ltd.（大阪大学 基礎工学部 卒業）\\
\textbf{ORCID：} 0009-0004-6753-4020\\
\textbf{作成日：} 2026年4月\\
\textbf{種別：} 研究ノート（幾何学的考察）\\
\textbf{DOI：} https://doi.org/10.5281/zenodo.19435162\\
\textbf{前論文：} - {[}1{]} Geometric Formulation of 4-Dimensional Spacetime via Gnomonic Projection. DOI: 10.5281/zenodo.19427780. - {[}2{]} 中心投影の幾何学的対称性：多軸モデルの数学的基盤. DOI: 10.5281/zenodo.19434932.

\begin{center}\rule{0.5\linewidth}{0.5pt}\end{center}

\subsection{概要}\label{ux6982ux8981}

前稿 {[}1{]} において，\(n\) 次元接超平面から \(S^n(R)\) への中心投影の誘導計量および Einstein テンソルを導出した．前稿 {[}2{]} において，この中心投影が5つの幾何学的対称性（離散安定性，軸の対等性，測地線偏差，主観座標系の変換可能性，大円上の向心加速度）を持つことを示し，特に\textbf{主観空間の内部観測者は背景空間の情報（座標の符号，曲率の方向，軸の絶対位置）を一切認識できない}ことを定理 5.2 として確立した．

本稿は，前2稿で確立された幾何学的枠組みに，第5の軸を追加する場合の構造的帰結を考察する．具体的には：

\begin{enumerate}
\def\labelenumi{\arabic{enumi}.}
\item
  \textbf{離散等方性の数論的要請：} 離散格子 \(\mathbb{Z}^{1+N}\) 上で等方的な光的測地線が存在するための条件は \(N\) の素因数分解に関する数論的制約（Fermat の二平方和定理）であり，\(N=3\) では破れ \(N=4\) で初めて成立する．これは離散時空が等方性を保つために第5の軸を要請することを示す．
\item
  \textbf{第5軸の量子化：} 結合定数 \(C\) と最小単位 \(e_0\) により定義される変換係数 \(Q = C/e_0\) のもとで，第5軸の主観座標 \(\tilde{\phi} = q/e_0\) は整数値をとり，コンパクト化を仮定することなく量子化が成立する．
\item
  \textbf{第5軸方向の曲率と \(1/r^2\) 構造：} 第5軸を曲率半径として中心投影を構成すると，主観空間における測地線偏差 \(|\xi|^2\) が距離の2乗に反比例する形で現れる．これは前稿 {[}2{]} 定理 5.2 (i) の帰結として導かれ，内部観測者の原理（埋め込み座標 \(Y^A\) は不可観測）と完全に整合する．
\end{enumerate}

本稿では物理的解釈は最小限にとどめ，幾何学的命題と前2稿からの厳密な帰結のみを述べる．主観座標と背景座標の混同，および埋め込み座標の物理的観測量としての扱いは，前2稿で禁止されている操作であり，本稿でも一貫して避ける．

\begin{center}\rule{0.5\linewidth}{0.5pt}\end{center}

\subsection{1. 前2稿からの前提}\label{ux524d2ux7a3fux304bux3089ux306eux524dux63d0}

\subsubsection{1.1 中心投影の基本構造（前稿 {[}1{]}）}\label{ux4e2dux5fc3ux6295ux5f71ux306eux57faux672cux69cbux9020ux524dux7a3f-1}

\(n\) 次元接超平面 \(\Pi_R\) から半径 \(R\) の超球面 \(S^n(R) \subset \mathbb{R}^{n+1}\) への中心投影 \(\Phi_A\) は，投影中心軸 \(A\) に対して

\[Y^\mu = \frac{Rx^\mu}{l_A} \quad (\mu \neq A), \qquad Y^A = \frac{R^2}{l_A}, \qquad l_A^2 = R^2 + \sum_{\mu \neq A} (x^\mu)^2 \tag{1.1}\]

で与えられる．引き戻し計量は

\[g_{\mu\nu} = \frac{R^2}{l_A^2}\left(\delta_{\mu\nu} - \frac{x_\mu x_\nu}{l_A^2}\right) \tag{1.2}\]

であり，\(n\) 次元 Einstein テンソルは

\[G_{\mu\nu} + \Lambda_n\, g_{\mu\nu} = 0, \qquad \Lambda_n = \frac{(n-1)(n-2)}{2R^2} \tag{1.3}\]

を満たす．

\subsubsection{1.2 5つの対称性（前稿 {[}2{]}）}\label{ux3064ux306eux5bfeux79f0ux6027ux524dux7a3f-2}

前稿 {[}2{]} で確立された5つの対称性を再掲する：

\begin{itemize}
\tightlist
\item
  \textbf{対称性 I（離散安定性）：} 任意の \(x^\mu \in \mathbb{R}\)（特に整数値）に対し \(l_A > 0\)，写像は正則．
\item
  \textbf{対称性 II（軸の対等性）：} 投影中心軸 \(A\) の選択に依らず，計量・曲率テンソルの構造は同一．
\item
  \textbf{対称性 III（測地線偏差）：} 測地線偏差 \(|\xi|^2 = g_{\mu\nu}\xi^\mu\xi^\nu\) は \([L^2]\) 次元のスカラー不変量であり，\textbf{座標の符号も曲率方向も内部観測者には弁別できない}．これが内部から測定しうる唯一の量である．
\item
  \textbf{対称性 IV（主観座標系の変換可能性）：} 異なる軸を中心とする主観座標系は \(C^\infty\) 級に相互変換可能．
\item
  \textbf{対称性 V（大円上の向心加速度）：} 主観空間の大円（測地線）に沿って等速円運動する内部観測者は，向心加速度 \(a = R\omega^2\) をスカラーとして測定する．方向は弁別できない．
\end{itemize}

\subsubsection{1.3 内部観測者の原理}\label{ux5185ux90e8ux89b3ux6e2cux8005ux306eux539fux7406}

前2稿を貫く基本原理は次の通りである：

\begin{quote}
\textbf{内部観測者の原理：} 主観空間 \(H_A^+\) の内部観測者は，引き戻し計量 \(g_{\mu\nu}\) とそれから導かれるスカラー量（測地線偏差，曲率スカラー等）のみを通じて測定を行う．埋め込み座標 \(Y^A\) や背景空間 \(\mathbb{R}^{n+1}\) の情報には直接アクセスできない．
\end{quote}

本稿のすべての主張は，この原理に従う．特に，加速度や力に対応する量を扱う場合は，\textbf{主観座標における測地線偏差}を経由して定義する．埋め込み座標 \(Y^A\) の時間微分を物理量として扱うことは，本原理に反するため行わない．

\begin{center}\rule{0.5\linewidth}{0.5pt}\end{center}

\subsection{2. 数論的動機：第5軸の必要性}\label{ux6570ux8ad6ux7684ux52d5ux6a5fux7b2c5ux8ef8ux306eux5fc5ux8981ux6027}

\subsubsection{2.1 離散格子上の等方的光的測地線}\label{ux96e2ux6563ux683cux5b50ux4e0aux306eux7b49ux65b9ux7684ux5149ux7684ux6e2cux5730ux7dda}

\textbf{観察 2.1（離散等方性の数論的条件）}\\
離散 Minkowski 格子 \(\mathbb{Z}^{1+N}\) 上で，光的等方測地線

\[\sum_{i=1}^{N} x_i^2 = t^2 \tag{2.1}\]

が非自明な整数解 \((x_1, \ldots, x_N, t) \in \mathbb{Z}^{N+1} \setminus \{0\}\) を持つ条件は，Fermat の二平方和定理（Lagrange の四平方和定理を含む）の系として次で与えられる：\(t^2\) が \(N\) 個の平方数の和として表せる必要がある．

\textbf{\(N = 3\) の場合：} \(t^2 = x_1^2 + x_2^2 + x_3^2\)．Legendre の三平方和定理により，\(t^2\) が \(4^a(8b+7)\) の形を含まないこと．特に等方的な解（\(x_1 = x_2 = x_3 = x\)）を要求すると \(t^2 = 3x^2\) となり，\(\sqrt{3} \notin \mathbb{Q}\) から非自明整数解は存在しない．

\textbf{\(N = 4\) の場合：} \(t^2 = x_1^2 + x_2^2 + x_3^2 + x_4^2\)．等方的な解（\(x_1 = \cdots = x_4 = x\)）を要求すると \(t^2 = 4x^2\)，すなわち \(t = 2x\)．これは任意の整数 \(x\) に対して整数解を与える．最小解は \((x_1, x_2, x_3, x_4, t) = (1,1,1,1,2)\)．\(\square\)

\textbf{命題 2.1（離散等方性のための次元要請）}\\
離散格子 \(\mathbb{Z}^{1+N}\) 上で等方的な光的測地線が存在するための最小の空間次元は \(N = 4\) である．

\textbf{証明} 観察 2.1 より \(N \leq 3\) では等方的整数解が存在せず，\(N = 4\) で初めて存在する．\(\square\)

\textbf{注意 2.1（命題 2.1 の解釈）}\\
命題 2.1 は，物理的な空間が4次元であることを主張するものではない．主張は次のように制限される：

\begin{quote}
もし離散格子上で「等方的な光的伝搬」を許容する幾何学的構造を要求するならば，必要な空間軸の数は最低4つである．
\end{quote}

3空間次元と1時間次元の枠組みに第4の空間軸を加える方法は数学的に複数考えうる．本稿では，この第4軸を\textbf{主観座標として量子化される第5の独立軸}と同定する立場を取る．この同定が物理的に妥当であるかは，本稿の範囲外である．

\subsubsection{2.2 第5軸の幾何学的位置づけ}\label{ux7b2c5ux8ef8ux306eux5e7eux4f55ux5b66ux7684ux4f4dux7f6eux3065ux3051}

前稿 {[}2{]} 対称性 II（軸の対等性）により，背景空間 \(\mathbb{R}^{n+1}\) のどの軸を投影中心に選んでも幾何学的構造は同一である．したがって既存の4軸 \((x, y, z, ct)\) に対して第5軸 \(\phi\) を追加することは，対称性 II の枠内で正当化される操作である．

ここで第5軸 \(\phi\) の主観座標としての性質は次のようにまとめられる：

\begin{itemize}
\tightlist
\item
  \(\phi\) は前稿 {[}2{]} 定理 4.1 の意味で他の4軸と対等な主観座標である
\item
  \(\phi\) の符号は前稿 {[}2{]} 定理 5.2 (ii) により内部観測者には弁別不可能
\item
  \(\phi\) の絶対値の2乗 \(\phi^2\) のみが計量を通じて内部観測者に寄与する
\end{itemize}

\begin{center}\rule{0.5\linewidth}{0.5pt}\end{center}

\subsection{3. 第5軸の離散化と量子化}\label{ux7b2c5ux8ef8ux306eux96e2ux6563ux5316ux3068ux91cfux5b50ux5316}

\subsubsection{3.1 結合定数による離散化}\label{ux7d50ux5408ux5b9aux6570ux306bux3088ux308bux96e2ux6563ux5316}

前稿 {[}2{]} で導入された結合定数 \(C\)（次元：長さ）を用いて，5軸すべてを離散化する：

\[\tilde{x}^\mu = \frac{x^\mu}{C} \in \mathbb{Z}, \quad \mu \in \{1,2,3,4,5\} \tag{3.1}\]

第4軸 \(\phi\) について，「最小単位」を \(e_0\)（次元：\(\phi\) と同じ）と書き，変換係数 \(Q\) を

\[Q = \frac{C}{e_0} \tag{3.2}\]

で定義する．元の物理量 \(q\)（次元：\(\phi/Q\)）に対して

\[\phi = Qq, \quad \tilde{\phi} = \frac{\phi}{C} = \frac{Qq}{C} = \frac{q}{e_0} \tag{3.3}\]

となり，離散化条件 \(\tilde{\phi} \in \mathbb{Z}\) から

\[q = n \cdot e_0, \quad n \in \mathbb{Z} \tag{3.4}\]

が得られる．

\subsubsection{3.2 量子化はコンパクト化を要さない}\label{ux91cfux5b50ux5316ux306fux30b3ux30f3ux30d1ux30afux30c8ux5316ux3092ux8981ux3055ux306aux3044}

\textbf{命題 3.1（離散格子からの量子化）}\\
離散化条件 \(\tilde{\phi} \in \mathbb{Z}\) のみから，第4軸の量子 \(q\) は最小単位 \(e_0\) の整数倍に限定される．第4軸方向のコンパクト化（周期境界条件）は仮定していない．

\textbf{証明} (3.3)--(3.4) は代数的恒等関係であり，仮定は \(\tilde{\phi} \in \mathbb{Z}\) のみである．第4軸の位相幾何学的構造（コンパクトか非コンパクトか）には言及していない．\(\square\)

\textbf{注意 3.1（Kaluza-Klein との対比）}\\
Kaluza-Klein 型の理論では量子化は第5次元のコンパクト化に伴う周期境界条件から導かれる．本稿では量子化は離散格子の整数性のみから導かれ，第4軸は他の軸と同じく非コンパクトでありうる．この差は本稿の枠組みが前稿 {[}2{]} 対称性 I（離散安定性）に基づいていることの直接的帰結である．

\textbf{注意 3.2（\(e_0\) の値について）}\\
本稿は \(e_0\) の数値を指定しない．\(e_0\) は本理論の自由パラメータではなく，何らかの外部基準（実験等）によって決定されるべき量である．本稿の幾何学的主張は \(e_0\) の具体的値に依らない．

\begin{center}\rule{0.5\linewidth}{0.5pt}\end{center}

\subsection{4. 5次元中心投影の構成}\label{ux6b21ux5143ux4e2dux5fc3ux6295ux5f71ux306eux69cbux6210}

\subsubsection{4.1 一般化された中心投影}\label{ux4e00ux822cux5316ux3055ux308cux305fux4e2dux5fc3ux6295ux5f71}

前稿 {[}1{]} §3 の構成を \(n=5\) に拡張する．背景空間 \(\mathbb{R}^6\) の符号系を

\[\epsilon_\mu = (+1, +1, +1, +1, -1, +1) \tag{4.1}\]

とする．ここで第6成分（投影中心軸方向）は前稿 {[}1{]} 付録C の符号統一公式に従い \(+1\) に固定する．第5成分 \(\epsilon_5 = -1\) は時間軸に対応する．

5次元接超平面 \(\Pi_R\)（座標 \((x^1, x^2, x^3, x^4, x^5)\)）から \(S^5(R) \subset \mathbb{R}^6\) への中心投影は

\[Y^\mu = \frac{Rx^\mu}{l} \quad (\mu = 1,\ldots,5), \quad Y^6 = \frac{R^2}{l} \tag{4.2}\]

\[l^2 = R^2 + \sum_{\mu=1}^{5} \epsilon_\mu (x^\mu)^2 = R^2 + (x^1)^2 + (x^2)^2 + (x^3)^2 + (x^4)^2 - (x^5)^2 \tag{4.3}\]

で与えられる．

\subsubsection{4.2 引き戻し計量}\label{ux5f15ux304dux623bux3057ux8a08ux91cf}

前稿 {[}1{]} 付録C の統一公式 (C.3) より，引き戻し計量は

\[g_{\mu\nu} = \frac{R^2}{l^2}\left(\epsilon_\mu \delta_{\mu\nu} - \frac{\epsilon_\mu \epsilon_\nu x_\mu x_\nu}{l^2}\right) \tag{4.4}\]

となる（\(\mu, \nu = 1,\ldots,5\)）．

\subsubsection{4.3 5次元 Einstein テンソル}\label{ux6b21ux5143-einstein-ux30c6ux30f3ux30bdux30eb}

前稿 {[}1{]} 注意 6.2 の一般式 \(\Lambda_n = (n-1)(n-2)/(2R^2)\) に \(n = 5\) を代入：

\[\Lambda_5 = \frac{(5-1)(5-2)}{2R^2} = \frac{6}{R^2} \tag{4.5}\]

\[G_{\mu\nu} + \frac{6}{R^2} g_{\mu\nu} = 0 \tag{4.6}\]

これは前稿 {[}1{]} §6 の主結果の \(n=5\) への直接拡張であり，新たな証明を要さない．

\subsubsection{4.4 離散版の正則性}\label{ux96e2ux6563ux7248ux306eux6b63ux5247ux6027}

前稿 {[}2{]} 対称性 I（離散安定性）の5次元への拡張．

\textbf{定理 4.1（5次元離散中心投影の正則性）}\\
\(\tilde{l}^2 = 1 + w^2 \tilde{\sigma}^2\)，\(\tilde{\sigma}^2 = \sum_{\mu=1}^{5} \epsilon_\mu (\tilde{x}^\mu)^2\)，\(w = C/R\) とする．

\textbf{(i) Euclidean 符号}（\(\epsilon_\mu = (+1,+1,+1,+1,+1)\)）：任意の整数座標 \(\tilde{x}^\mu \in \mathbb{Z}\) に対し \(\tilde{\sigma}^2 \geq 0\)，したがって \(\tilde{l}^2 \geq 1 > 0\)．大域的に正則．

\textbf{(ii) Lorentzian 符号}（\(\epsilon_\mu = (+1,+1,+1,+1,-1)\)）：\(\tilde{\sigma}^2 = (\tilde{x}^1)^2 + (\tilde{x}^2)^2 + (\tilde{x}^3)^2 + (\tilde{x}^4)^2 - (\tilde{x}^5)^2\)．\(\tilde{l}^2 > 0\) となる領域は

\[1 + w^2\left[(\tilde{x}^1)^2 + (\tilde{x}^2)^2 + (\tilde{x}^3)^2 + (\tilde{x}^4)^2\right] > w^2 (\tilde{x}^5)^2 \tag{4.7}\]

で与えられる．

\textbf{証明} 前稿 {[}2{]} §3 定理 3.1 の証明を5次元に適用する．Euclidean 版では \(\tilde{\sigma}^2 \geq 0\) が各項の非負性から自明に従う．Lorentzian 版では (4.7) が \(\tilde{l}^2 > 0\) の必要十分条件として代数的に直接得られる．\(\square\)

\begin{center}\rule{0.5\linewidth}{0.5pt}\end{center}

\subsection{\texorpdfstring{5. 第5軸方向の中心投影と \(1/r^2\) 構造}{5. 第5軸方向の中心投影と 1/r\^{}2 構造}}\label{ux7b2c5ux8ef8ux65b9ux5411ux306eux4e2dux5fc3ux6295ux5f71ux3068-1r2-ux69cbux9020}

本節では，前稿 {[}2{]} 対称性 II（軸の対等性）を利用して，第4軸 \(\phi\)（あるいは符号系を入れ替えた別の軸）を投影中心とする中心投影を構成し，主観空間における測地線偏差の振る舞いを調べる．

\subsubsection{5.1 軸の選択と背景の設定}\label{ux8ef8ux306eux9078ux629eux3068ux80ccux666fux306eux8a2dux5b9a}

前稿 {[}2{]} 対称性 II により，5軸 \((x^1, x^2, x^3, x^4, x^5) = (x, y, z, \phi, ct)\) のうちどの軸を投影中心軸として選んでも幾何学的構造は同一である．ここでは時間軸 \(ct\) ではなく第4軸 \(\phi\) を投影中心軸とした構成を考える．

投影中心軸を第4軸 \(\phi\) に取ると，主観座標は残りの4軸 \((x, y, z, ct)\) となり，曲率半径は

\[\phi_0 = \text{第4軸方向の曲率半径} \tag{5.1}\]

である．接超平面は4次元 Lorentzian であり，4次元主観座標 \((x, y, z, ct)\) を持つ．

このときの引き戻し計量は，前稿 {[}1{]} 付録C 統一公式 (C.3) において \(A = 4\) とした場合に相当し：

\[g_{\mu\nu} = \frac{\phi_0^2}{l_\phi^2}\left(\eta_{\mu\nu} - \frac{\eta_{\mu\alpha}\eta_{\nu\beta} x^\alpha x^\beta}{l_\phi^2}\right), \quad l_\phi^2 = \phi_0^2 + \sigma^2 \tag{5.2}\]

ここで \(\eta_{\mu\nu} = \text{diag}(+1,+1,+1,-1)\)，\(\sigma^2 = x^2 + y^2 + z^2 - c^2 t^2\)．

\subsubsection{5.2 主観空間の測地線偏差}\label{ux4e3bux89b3ux7a7aux9593ux306eux6e2cux5730ux7ddaux504fux5dee}

前稿 {[}2{]} §5 で確立された測地線偏差 \(|\xi|^2 = g_{\mu\nu}\xi^\mu\xi^\nu\) を，計量 (5.2) のもとで評価する．

\textbf{命題 5.1（測地線偏差の距離依存性）}\\
計量 (5.2) のもとで，主観空間内の2点を結ぶ短距離測地線族の偏差ベクトル \(\xi^\mu\) に対し，測地線偏差 \(|\xi|^2\) は曲率半径 \(\phi_0\) と空間距離 \(r = \sqrt{x^2+y^2+z^2}\) について次のスケーリングを持つ：

\[|\xi|^2 \sim \frac{\phi_0^2}{(\phi_0^2 + \sigma^2)^2} \cdot |\xi_0|^2 \tag{5.3}\]

ここで \(|\xi_0|^2\) は接点近傍での偏差の参照値である．特に静的な空間断面（\(ct = 0\)）における大距離極限 \(r \gg \phi_0\) では：

\[|\xi|^2 \sim \frac{\phi_0^2}{r^4} \cdot |\xi_0|^2 \tag{5.4}\]

\textbf{証明} 計量 (5.2) の前因子は \(\phi_0^2/l_\phi^2\) であり，\(|\xi|^2 = g_{\mu\nu}\xi^\mu\xi^\nu\) の各項はこの前因子に比例する．第二項 \(-\eta_{\mu\alpha}\eta_{\nu\beta}x^\alpha x^\beta/l_\phi^2\) は補正項であり，\(r \gg \phi_0\) では支配的にならない場合，\(|\xi|^2\) の主要なスケーリングは前因子で決まる．具体的には

\[|\xi|^2 \approx \frac{\phi_0^2}{l_\phi^2}\eta_{\mu\nu}\xi^\mu\xi^\nu = \frac{\phi_0^2}{\phi_0^2 + \sigma^2}\eta_{\mu\nu}\xi^\mu\xi^\nu\]

参照値 \(|\xi_0|^2 = \eta_{\mu\nu}\xi^\mu\xi^\nu\) を分離すると (5.3) を得る．静的断面 \(ct = 0\) で \(\sigma^2 = r^2\)，\(r \gg \phi_0\) で \(l_\phi^2 \approx r^2\)，したがって \(|\xi|^2 \sim \phi_0^2/r^4 \cdot |\xi_0|^2\)．\(\square\)

\textbf{注意 5.1（主観座標における意味）}\\
命題 5.1 は測地線偏差 \(|\xi|^2\) という前稿 {[}2{]} 定理 5.2 で内部観測者に測定可能と確立された量についての記述である．埋め込み座標 \(Y^\mu\) には一切言及していない．主観空間の内部観測者は，測地線偏差の距離依存性 \(|\xi|^2 \propto \phi_0^2/r^4\) を測定することができ，このスケーリングが内部観測者に見える \(\phi_0\) の効果である．

\textbf{注意 5.2（弁別不可能性）}\\
前稿 {[}2{]} 定理 5.2 (ii) により，\(|\xi|^2\) はスカラー不変量であるため，内部観測者は座標の符号も曲率方向も弁別できない．したがって命題 5.1 の \(|\xi|^2 \propto \phi_0^2/r^4\) は \(\phi_0\) の符号にも \(r\) の方向にも依存しない正の量である．「曲率方向」や「正負の電荷」といった概念は背景空間 \(\mathbb{R}^6\) においてのみ定義され，主観空間の内在幾何には現れない．

\subsubsection{5.3 結果の整理}\label{ux7d50ux679cux306eux6574ux7406}

本稿で導かれた幾何学的結果は次の通りである：

\begin{enumerate}
\def\labelenumi{\arabic{enumi}.}
\tightlist
\item
  第5軸 \(\phi\) の追加は前稿 {[}2{]} 対称性 II により正当化される
\item
  第5軸の離散化条件 \(\tilde{\phi} \in \mathbb{Z}\) から量子化 \(q = n e_0\) が代数的に従う（命題 3.1）
\item
  第5軸を投影中心軸とした中心投影の引き戻し計量から，主観空間の測地線偏差 \(|\xi|^2\) は \(\phi_0^2/r^4\) のスケーリングを持つ（命題 5.1）
\item
  内部観測者は \(|\xi|^2\) のスカラー値のみを測定でき，\(\phi_0\) や \(r\) の方向や符号は弁別できない（前稿 {[}2{]} 定理 5.2 の帰結）
\end{enumerate}

これらはすべて主観座標と内在計量のみで記述された幾何学的命題であり，埋め込み座標 \(Y^A\) の物理的観測量としての扱いを一切含まない．内部観測者の原理（§1.3）は本稿全体を通じて厳守されている．

\begin{center}\rule{0.5\linewidth}{0.5pt}\end{center}

\subsection{6. 物理的解釈に関する留保}\label{ux7269ux7406ux7684ux89e3ux91c8ux306bux95a2ux3059ux308bux7559ux4fdd}

\subsubsection{6.1 本稿が主張しないこと}\label{ux672cux7a3fux304cux4e3bux5f35ux3057ux306aux3044ux3053ux3068}

本稿は次のことを\textbf{主張しない}：

\begin{itemize}
\tightlist
\item
  第4軸 \(\phi\) が物理的な「電荷」次元と同一であること
\item
  命題 5.1 のスケーリング \(|\xi|^2 \propto \phi_0^2/r^4\) が Coulomb の法則と同型であること
\item
  内部観測者が感じる「加速度」が命題 5.1 から導かれること
\item
  Maxwell 方程式が本稿の枠組みから導出されること
\item
  微細構造定数 \(\alpha\) や電子電荷 \(e\) の値が本稿から決定されること
\end{itemize}

これらの物理的同定はいずれも本稿の幾何学的命題を超える追加の仮定を要する．特に「測地線偏差を加速度として読み替える」操作は，計量と粒子の運動方程式の関係を別途定式化する必要があり，本稿の範囲外である．

\subsubsection{6.2 本稿が主張すること}\label{ux672cux7a3fux304cux4e3bux5f35ux3059ux308bux3053ux3068}

本稿が主張するのは次の幾何学的命題のみである：

\begin{itemize}
\tightlist
\item
  離散格子上の等方光的測地線は第5軸を要請する（命題 2.1）
\item
  第5軸の離散化から代数的に量子化が従う（命題 3.1）
\item
  第5軸を投影中心軸とした中心投影の主観空間における測地線偏差は \(\phi_0^2/r^4\) にスケールする（命題 5.1）
\end{itemize}

これらはすべて前2稿で確立された幾何学的枠組みの直接的帰結であり，本稿は新たな公理や仮定を導入していない．

\subsubsection{6.3 今後の課題}\label{ux4ecaux5f8cux306eux8ab2ux984c}

本稿の幾何学的結果と既存の物理理論との対応関係は，独立した研究課題である．特に：

\begin{itemize}
\tightlist
\item
  命題 5.1 のスケーリングと粒子の運動方程式（測地線方程式）との関係の定式化
\item
  第5軸の量子 \(e_0\) の物理的同定
\item
  4次元時空への有効理論としての射影と，その整合性の確認
\item
  Lorentz 不変性（前稿 {[}2{]} で確立）と本稿の第5軸構造の整合的接続
\end{itemize}

これらは本稿の幾何学的基盤の上に構築されるべき次段階の研究である．

\begin{center}\rule{0.5\linewidth}{0.5pt}\end{center}

\subsection{結果の整理}\label{ux7d50ux679cux306eux6574ux7406-1}

{\def\LTcaptype{none} % do not increment counter
\begin{longtable}[]{@{}
  >{\raggedright\arraybackslash}p{(\linewidth - 6\tabcolsep) * \real{0.2069}}
  >{\raggedright\arraybackslash}p{(\linewidth - 6\tabcolsep) * \real{0.5172}}
  >{\raggedright\arraybackslash}p{(\linewidth - 6\tabcolsep) * \real{0.0690}}
  >{\raggedright\arraybackslash}p{(\linewidth - 6\tabcolsep) * \real{0.2069}}@{}}
\toprule\noalign{}
\begin{minipage}[b]{\linewidth}\raggedright
項目
\end{minipage} & \begin{minipage}[b]{\linewidth}\raggedright
結果
\end{minipage} & \begin{minipage}[b]{\linewidth}\raggedright
式番号
\end{minipage} & \begin{minipage}[b]{\linewidth}\raggedright
根拠
\end{minipage} \\
\midrule\noalign{}
\endhead
\bottomrule\noalign{}
\endlastfoot
数論的動機 & \(N=4\) で初めて等方光的整数解 \((1,1,1,1,2)\) が存在 & (2.1) & 命題 2.1 \\
第5軸の正当化 & 軸の対等性により追加可能 & §2.2 & 前稿 {[}2{]} 対称性 II \\
量子化 & \(\tilde{\phi} \in \mathbb{Z}\) から \(q = ne_0\) が代数的に従う & (3.4) & 命題 3.1 \\
コンパクト化不要 & 量子化は離散格子の整数性のみに依存 & 注意 3.1 & 命題 3.1 \\
5次元計量 & 前稿 {[}1{]} 付録C の統一公式で記述 & (4.4) & 前稿 {[}1{]} \\
5次元 \(\Lambda\) & \(\Lambda_5 = 6/R^2\) & (4.5) & 前稿 {[}1{]} 注意 6.2 \\
5次元正則性 & Euclidean 全域，Lorentzian は (4.7) 内 & (4.7) & 定理 4.1 \\
測地線偏差のスケーリング & \$ & \xi & \^{}2 \propto \phi\_0\textsuperscript{2/r}4\$ \\
内部観測者の弁別不可能性 & \(\phi_0\) の符号も \(r\) の方向も不可弁別 & 注意 5.2 & 前稿 {[}2{]} 定理 5.2 (ii) \\
\end{longtable}
}

\begin{center}\rule{0.5\linewidth}{0.5pt}\end{center}

\subsection{参考文献}\label{ux53c2ux8003ux6587ux732e}

{[}1{]} Kihara, N. (2026). \emph{Geometric Formulation of 4-Dimensional Spacetime via Gnomonic Projection}. DOI: 10.5281/zenodo.19427780.

{[}2{]} Kihara, N. (2026). \emph{中心投影の幾何学的対称性：多軸モデルの数学的基盤}. DOI: 10.5281/zenodo.19434932.

{[}3{]} Hardy, G. H. and Wright, E. M., \emph{An Introduction to the Theory of Numbers}, 6th ed., Oxford University Press, 2008, Chapter 20 (Lagrange's four-square theorem, Legendre's three-square theorem).

\begin{center}\rule{0.5\linewidth}{0.5pt}\end{center}

\emph{本稿は前2稿 {[}1, 2{]} で確立された幾何学的枠組みの第5軸への拡張を扱う研究ノートである．物理的解釈は §6 に記した範囲を超えない．内部観測者の原理を一貫して厳守し，埋め込み座標の物理的観測量としての扱いは行わない．}

\end{document}
